\documentclass[english,11pt]{letter}
\usepackage[T1]{fontenc}
\usepackage[latin1]{inputenc}
\usepackage{babel}
\usepackage[top=.5in,bottom=.5in,left=1in,right=1in,noheadfoot]{geometry}
\usepackage{graphicx}
\usepackage{url}

\address{\sffamily
  \begin{tabular}{rr}
    Prof.~Dr.~Ant\^{o}nio Augusto Fr\"{o}hlich & \raisebox{-0.9cm}[0pt][0pt]{\resizebox*{!}{1.2cm}{\includegraphics{lisha}}}\\
    Federal University of Santa Catarina &\\
    Software/Hardware Integration Laboratory &\\
    \multicolumn{2}{r}{PO Box 476, 88040-900 Florian\'{o}polis, SC, Brazil}\\
    \multicolumn{2}{r}{Phone: +55 48 3721-9516}\\
    \multicolumn{2}{r}{E-Mail: guto@lisha.ufsc.br}\\
    \multicolumn{2}{r}{WWW: http://lisha.ufsc.br/Guto}\\
    \\
    \multicolumn{2}{r}{\today}
  \end{tabular}
}

\date{}

\signature{
  \begin{center}
   %\raisebox{-.9cm}[0pt][0pt]{\includegraphics[scale=1]{/home/guto/personal/signature}}\\
    \hrulefill\\
    Ant\^{o}nio Augusto Fr\"{o}hlich
  \end{center}
}

\begin{document}

\begin{letter}{ELSEVIER Computers and Electrical Engineering Journal\\
    Editor-in-Chief\\
    Professor Dr. Manu Malek}
\opening{Dear Prof. Malek,}
\setlength{\parskip}{1.5ex}

Thank you very much for the review of the manuscript \textbf{Periodic
  Timers Revisited: the Real-time Embedded System Perspective}. A
revised version that incorporates all your remarks and also the
reviewers' suggestions is being submitted online. Graphical abstract and
research highlights have been added. The addressing of reviewers'
comments is discussed bellow:

\begin{itemize}
\item \textbf{Precise description of periodic and single-shot
    mechanisms} 

  We have extended the explanation about periodic and single-shot timers
  in the introduction and we also included a figure to depict the
  described example.
	
\item \textbf{Paper's motivation and results}
	
  We made both motivation and results clearer in the introduction. Our
  main purpose is to weave some of the prejudice against periodic
  timers, which actually owe much of their bad reputation to poor
  implementations. Besides, we present concrete information about the
  limitations of single-shot timers that might be specially useful for
  those endeavoring new real-time support systems.
	
\item \textbf{Improving related work section}
	
  We have fully rewritten the related work section, including additional
  timing solutions from other groups and describing each one in more
  detail. Well-known Real-Time Operating Systems (e.g. VxWorks, QNX, RTAI,
  RTLinux) are now featured.
	
\item \textbf{Motivation for EPOS}
	
  As we stated in the final remarks of Section 4, a smart implementation
  of a time management mechanism does not depend on the real-time
  operating system. We used EPOS in our experiments because it features
  an elegant and clean implementation that makes additions and
  extensions rather easier than for traditional, non component-based
  systems. An open-source version of EPOS is available online at
  \url{http://epos.lisha.ufsc.br}.
	
\item \textbf{Conclusions}
	
  We clearly cited the case where a periodic timer outperforms an
  equivalent single-shot timer (when requested periods exceed the
  maximum hardware period) and also explained why this happens (overhead
  of reprogramming the hardware timer and software tick counting).
  Indeed, we added the result of the multi-threaded environment
  experiment, in which the periodic interrupt handler presented better
  execution timer up to 5 threads in comparison to the single-shot
  interrupt handler.
	
\item \textbf{Figures placement}
	
  We made a complete review in the figures placement. We also changed
  the footnote on page 2 by a Figure as recommended by the reviewer \#2.
	
\item \textbf{Tables 1 and 2}
	
  We changed the word targeted by expected in Tables 1 and 2 in order to
  improve the comprehension.
	
\item \textbf{Abbreviations}
	
  We have added a reference when AVR is used by the first time. AVR is
  not an abbreviation, but the name that Atmel gave for its
  architecture. We also made a review to spell out each acronym in the
  first time it is used.
	
\end{itemize}

	
Looking forward to have the manuscript published,

\closing{~}

\end{letter}

\end{document}
